% Einheit 1 – Nr. 13–29
\setcounter{aufgabe}{12}
\einheitenkopf[e1][1 Linear oder exponentiell]{Einheit 1 von 5 · Lineares und exponentielles Wachstum unterscheiden}
\zweigzeile{Hier lernst du, lineares und exponentielles Wachstum an Tabelle, Text und Graph zu unterscheiden · neu in diesem Jahr · P10 oft · baut auf: Punkte eintragen, gleicher Betrag je Schritt, Erhöhen um p\,\%}

\begin{aufgabe}{Ich erkenne, ob jedes Mal derselbe Betrag oder derselbe Prozentsatz dazukommt.}
Kreuze an. Rechne nichts.
\begin{teile}
\teil Ein Sparschwein bekommt jeden Monat 15\,€ dazu. \\ \kreuz{gleicher Betrag}\kreuz{gleicher Prozentsatz}
\teil Ein Guthaben wächst jedes Jahr um 2\,\%. \\ \kreuz{gleicher Betrag}\kreuz{gleicher Prozentsatz}
\teil Eine Bohnenpflanze wird jeden Tag 3\,cm höher. \\ \kreuz{gleicher Betrag}\kreuz{gleicher Prozentsatz}
\teil Die Zahl der Nutzer einer App steigt jeden Monat um 12\,\%. \\ \kreuz{gleicher Betrag}\kreuz{gleicher Prozentsatz}
\teil Ein Auto verliert jedes Jahr 15\,\% seines Werts. \\ \kreuz{gleicher Betrag}\kreuz{gleicher Prozentsatz}
\end{teile}
\end{aufgabe}

\begin{aufgabe}{Ich erkenne, ob in einer Tabelle die Differenz oder der Quotient gleich bleibt.}
Unter den Werten stehen schon die Differenzen und die Quotienten benachbarter Werte. Kreuze an, was gleich bleibt.
\begin{teile}
\teil Werte: 4; \ 12; \ 36; \ 108 \quad Differenzen: 8; \ 24; \ 72 \quad Quotienten: 3; \ 3; \ 3 \\ \kreuz{die Differenz}\kreuz{der Quotient}
\teil Werte: 50; \ 65; \ 80; \ 95 \quad Differenzen: 15; \ 15; \ 15 \quad Quotienten: 1,3; \ $\approx 1{,}23$; \ $\approx 1{,}19$ \\ \kreuz{die Differenz}\kreuz{der Quotient}
\teil Werte: 400; \ 200; \ 100; \ 50 \quad Differenzen: $-200$; \ $-100$; \ $-50$ \quad Quotienten: 0,5; \ 0,5; \ 0,5 \\ \kreuz{die Differenz}\kreuz{der Quotient}
\teil Werte: 90; \ 81; \ 72; \ 63 \quad Differenzen: $-9$; \ $-9$; \ $-9$ \quad Quotienten: 0,9; \ $\approx 0{,}89$; \ $\approx 0{,}88$ \\ \kreuz{die Differenz}\kreuz{der Quotient}
\end{teile}
\end{aufgabe}

\begin{aufgabe}{Ich erkenne, wo ein Graph die senkrechte Achse trifft.}
Markiere den Punkt, in dem der Graph die senkrechte Achse trifft. Rechne nichts.

\begin{ksys}[xmin=0,xmax=5,ymin=0,ymax=8,ystep=2,ablesen]
\gerade{1.2}{0}{A}
\end{ksys}\ksysabstand\begin{ksys}[xmin=0,xmax=5,ymin=0,ymax=8,ystep=2,ablesen]
\funktion{2*1.35^\x}{B}
\end{ksys}\ksysabstand\begin{ksys}[xmin=0,xmax=5,ymin=0,ymax=8,ystep=2,ablesen]
\gerade{0.8}{3}{C}
\end{ksys}

\begin{teile}
\teil bei Graph A
\teil bei Graph B
\teil bei Graph C
\teil Welcher Graph beginnt im Ursprung? \quad \kreuz{A}\kreuz{B}\kreuz{C}
\end{teile}
\end{aufgabe}

\verfahren{Wachstumsart an Tabelle und Text erkennen}

\begin{aufgabe}{Ich kann prüfen, ob in einer Tabelle immer derselbe Betrag dazukommt.}
Bilde die Differenzen benachbarter Werte. Kreuze an, ob sie gleich sind.
\begin{teile}
\teil 3; \ 7; \ 11; \ 15 \quad Differenzen: \leerfeld \quad gleich? \janein
\teil 2; \ 4; \ 8; \ 16 \quad Differenzen: \leerfeld \quad gleich? \janein
\teil 10; \ 25; \ 40; \ 55 \quad Differenzen: \leerfeld \quad gleich? \janein
\teil 5; \ 10; \ 20; \ 40 \quad Differenzen: \leerfeld \quad gleich? \janein
\teil 1,2; \ 1,8; \ 2,4; \ 3,0 \quad Differenzen: \leerfeld \quad gleich? \janein
\teil 90; \ 72; \ 54; \ 36 \quad Differenzen: \leerfeld \quad gleich? \janein
\end{teile}
\end{aufgabe}

\begin{aufgabe}{Ich kann prüfen, ob in einer Tabelle immer mit derselben Zahl multipliziert wird.}
Bilde die Quotienten: jeden Wert geteilt durch den Wert davor. Kreuze an, ob sie gleich sind.
\begin{teile}
\teil 3; \ 6; \ 12; \ 24 \quad Quotienten: \leerfeld \quad gleich? \janein
\teil 5; \ 15; \ 45; \ 135 \quad Quotienten: \leerfeld \quad gleich? \janein
\teil 2; \ 4; \ 6; \ 8 \quad Quotienten: \leerfeld \quad gleich? \janein
\teil 40; \ 50; \ 62,5; \ 78,125 \quad Quotienten: \leerfeld \quad gleich? \janein
\teil 800; \ 400; \ 200; \ 100 \quad Quotienten: \leerfeld \quad gleich? \janein
\teil 100; \ 110; \ 120; \ 130 \quad Quotienten: \leerfeld \quad gleich? \janein
\end{teile}
\end{aufgabe}

\begin{aufgabe}{Ich kann an einer Tabelle entscheiden, ob ein Wert linear oder exponentiell wächst.}
Entscheide und begründe in einem Satz.
\begin{teile}
\teil Kontostand nach 0, 1, 2, 3 Jahren: 200\,€; \ 260\,€; \ 320\,€; \ 380\,€
\teil Zahl der Follower nach 0, 1, 2, 3 Monaten: 500; \ 550; \ 605; \ 665,5
\teil Fläche eines Pilzgeflechts nach 0, 1, 2, 3 Wochen: 64\,cm²; \ 96\,cm²; \ 144\,cm²; \ 216\,cm²
\end{teile}
\end{aufgabe}

\begin{aufgabe}{Ich kann aus einem Text entscheiden, ob ein Wert linear oder exponentiell zu- oder abnimmt.}
Entscheide und begründe in einem Satz.
\begin{teile}
\teil Ein junger Baum wird jedes Jahr 40\,cm höher.
\teil Die Einwohnerzahl einer Stadt wächst jedes Jahr um 1,5\,\%.
\teil Ein Handy verliert jedes Jahr 25\,\% seines Werts.
\teil Aus einem Tank laufen jede Minute 12 Liter ab.
\teil Ein Fahrradverleih hatte diese Umsätze:

\sachtabelle{lcccc}{Jahr & 2021 & 2022 & 2023 & 2024}{Umsatz in € & 240\,000 & 252\,000 & 264\,600 & 277\,830}

Die Inhaberin sagt: „Der Umsatz wächst linear, denn er steigt jedes Jahr um 5\,\%.“ Entscheide, ob sie recht hat, und begründe. (P10 2018 OS)
\end{teile}
\end{aufgabe}

\verfahren{Graph zu einem Wachstum auswählen}

\begin{aufgabe}{Ich erkenne, ob ein Graph eine Gerade oder eine gekrümmte Kurve ist.}
Kreuze an.

\begin{ksys}[xmin=0,xmax=5,ymin=0,ymax=8,ystep=2,ablesen]
\gerade{1}{1}{A}
\end{ksys}\ksysabstand\begin{ksys}[xmin=0,xmax=5,ymin=0,ymax=8,ystep=2,ablesen]
\funktion{1.5^\x}{B}
\end{ksys}

\begin{ksys}[xmin=0,xmax=5,ymin=0,ymax=8,ystep=2,ablesen]
\funktion{0.5*2^\x}{C}
\end{ksys}\ksysabstand\begin{ksys}[xmin=0,xmax=5,ymin=0,ymax=8,ystep=2,ablesen]
\gerade{0.6}{2}{D}
\end{ksys}

\begin{teilezwei}
\tz A: \kreuz{Gerade}\kreuz{Kurve} & \tz B: \kreuz{Gerade}\kreuz{Kurve} \\
\tz C: \kreuz{Gerade}\kreuz{Kurve} & \tz D: \kreuz{Gerade}\kreuz{Kurve} \\
\end{teilezwei}
\end{aufgabe}

\begin{aufgabe}{Ich kann den passenden Graphen zu einem Wachstum auswählen.}
Auf der senkrechten Achse steht der Bestand in Tausend, auf der waagerechten die Zeit in Jahren. Kreuze an.

\begin{ksys}[xmin=0,xmax=5,ymin=0,ymax=8,ystep=2,ablesen]
\gerade{1.2}{0}{A}
\end{ksys}\ksysabstand\begin{ksys}[xmin=0,xmax=5,ymin=0,ymax=8,ystep=2,ablesen]
\funktion{2*1.4^\x}{B}
\end{ksys}\ksysabstand\begin{ksys}[xmin=0,xmax=5,ymin=0,ymax=8,ystep=2,ablesen]
\gerade{1.2}{2}{C}
\end{ksys}

\begin{teile}
\teil Ein Bestand von 2000 Tieren wächst jedes Jahr um 40\,\%. Welcher Graph kommt nicht in Frage, weil er nicht beim Anfangswert beginnt? \quad \kreuz{A}\kreuz{B}\kreuz{C}
\teil Welcher Graph passt zu diesem Bestand? \quad \kreuz{A}\kreuz{B}\kreuz{C}
\end{teile}

\begin{ksys}[xmin=0,xmax=5,ymin=0,ymax=8,ystep=2,ablesen]
\gerade{-1.4}{8}{D}
\end{ksys}\ksysabstand\begin{ksys}[xmin=0,xmax=5,ymin=0,ymax=8,ystep=2,ablesen]
\funktion{8*0.6^\x}{}
\node[fill=white,inner sep=1pt] at (1.7,5.4) {$E$};
\end{ksys}\ksysabstand\begin{ksys}[xmin=0,xmax=5,ymin=0,ymax=8,ystep=2,ablesen]
\funktion{1.5^\x}{F}
\end{ksys}

\begin{teile}
\teil Ein Bestand von 8000 nimmt jedes Jahr um 40\,\% ab. \quad \kreuz{D}\kreuz{E}\kreuz{F}
\teil Ein Bestand von 8000 nimmt jedes Jahr um 1400 ab. \quad \kreuz{D}\kreuz{E}\kreuz{F}
\end{teile}
\end{aufgabe}

\begin{aufgabe}{Ich kann begründen, warum ein Graph nicht zu einem Wachstum passt.}
Die Zahl der Nutzer einer App beginnt bei 3000 und wächst jeden Monat um 25\,\%. Senkrechte Achse: Nutzer in Tausend, waagerechte Achse: Zeit in Monaten. Begründe in einem Satz.

\begin{ksys}[xmin=0,xmax=5,ymin=0,ymax=8,ystep=2,ablesen]
\gerade{0.9}{3}{G}
\end{ksys}\ksysabstand\begin{ksys}[xmin=0,xmax=5,ymin=0,ymax=8,ystep=2,ablesen]
\funktion{0.35*\x^2}{H}
\end{ksys}\ksysabstand\begin{ksys}[xmin=0,xmax=5,ymin=0,ymax=8,ystep=2,ablesen]
\funktion{3*1.25^\x}{I}
\end{ksys}

\begin{teile}
\teil Warum passt der Graph G nicht?
\teil Warum passt der Graph H nicht?
\end{teile}

In einem See leben 6000 Fische. Wegen einer Krankheit nimmt ihre Zahl jedes Jahr um 30\,\% ab. Senkrechte Achse: Fische in Tausend, waagerechte Achse: Zeit in Jahren.

\begin{ksys}[xmin=0,xmax=5,ymin=0,ymax=8,ystep=2,ablesen]
\funktion{6*0.7^\x}{}
\node[fill=white,inner sep=1pt] at (1.7,5) {$J$};
\end{ksys}\ksysabstand\begin{ksys}[xmin=0,xmax=5,ymin=0,ymax=8,ystep=2,ablesen]
\gerade{-1.2}{6}{K}
\end{ksys}\ksysabstand\begin{ksys}[xmin=0,xmax=5,ymin=0,ymax=8,ystep=2,ablesen]
\funktion{8*0.7^\x}{L}
\end{ksys}

\begin{teile}
\teil Gib an, welcher Graph passt, und begründe für jeden der beiden anderen Graphen, warum er nicht passt. (P10 2026 FOR)
\end{teile}
\end{aufgabe}

\verfahren{Wertepaare als Punkte darstellen}

\begin{aufgabe}{Ich kann Wertepaare eines Wachstums als Punkte eintragen.}
Eine Zellkultur wird beobachtet ($t$: Zeit in Stunden, $N$: Zahl der Zellen). Trage die Punkte aus der Tabelle ein.

\wertetabelle[100,200,400]{t}{N}{0,1,2}

\begin{teile}
\teil Punkt für $t = 0$
\teil Punkt für $t = 1$
\teil Punkt für $t = 2$
\end{teile}

\begin{ksys}[xmin=0,xmax=4,ymin=0,ymax=500,ystep=100,xlabel=$t$ in h,ylabel=$N$]
\end{ksys}

\end{aufgabe}

\begin{aufgabe}{Ich kann Wertepaare eines Wachstums als Punkte eintragen – weiter.}
Eine zweite Kultur ($t$ in Stunden, $N$ Zahl der Zellen). Die ersten beiden Punkte sind schon eingetragen.

\wertetabelle[100,150,225,{337{,}5}]{t}{N}{0,1,2,3}

\begin{teile}
\teil Trage den Punkt für $t = 2$ genau ein.
\teil Trage den Punkt für $t = 3$ genau ein.
\teil Verbinde die vier Punkte zu einer Kurve.
\end{teile}

\begin{ksys}[xmin=0,xmax=4,ymin=0,ymax=400,ystep=50,xlabel=$t$ in h,ylabel=$N$]
\punkt{0}{100}{}
\punkt{1}{150}{}
\end{ksys}

\end{aufgabe}

\begin{aufgabe}{Ich kann für Wertepaare eine passende Achseneinteilung wählen.}
Ein E-Bike verliert jedes Jahr 20\,\% seines Werts.

\sachtabelle{lccccc}{Zeit in Jahren & 0 & 1 & 2 & 3 & 4}{Wert in € & 2000 & 1600 & 1280 & 1024 & 819,20}

\begin{teile}
\teil Teile beide Achsen passend ein und beschrifte sie.
\teil Trage die fünf Wertepaare als Punkte ein.
\teil Verbinde die Punkte zu einer Kurve. (P10 2017 OS)
\end{teile}

\leeresgitter{10}{10}

\end{aufgabe}

\begin{aufgabe}{Ich finde den Fehler bei der Entscheidung zwischen linear und exponentiell.}
Finde den Fehler, benenne ihn und korrigiere die Aussage.
\begin{teile}
\teil Ein Verein hat nacheinander 300, 330, 363 und 399,3 Mitglieder. Mia sagt: „Das Wachstum ist linear, denn jedes Jahr kommen 10\,\% dazu.“
\teil Ein Guthaben beginnt bei 500\,€ und wächst jedes Jahr um 2\,\%. Ben zeichnet dazu eine Gerade durch den Ursprung, „weil jedes Jahr gleich viel Prozent dazukommen“.
\end{teile}
\end{aufgabe}

\begin{aufgabe}{Ich kann begründen, warum bei gleichem Prozentsatz der Zuwachs immer größer wird.}
Begründe.
\begin{teile}
\teil Ein Guthaben von 1000\,€ wächst jedes Jahr um 10\,\%. Zeige an den ersten beiden Jahren, dass der Zuwachs in Euro größer wird.
\teil Warum heißt „jedes Jahr gleich viel Prozent“ nicht „jedes Jahr gleich viel Euro“?
\end{teile}
\end{aufgabe}

\begin{aufgabe}{Ich kann Grenzen eines Wachstumsmodells nennen.}
Begründe.
\begin{teile}
\teil Eine Sonnenblume ist 25\,cm hoch und wächst in den ersten Wochen jede Woche um 40\,\%. Kann sie ein ganzes Jahr so weiterwachsen?
\teil In einer Schule mit 800 Schülern verdoppelt sich jeden Tag die Zahl derer, die ein Gerücht kennen. Warum kann das nicht lange so weitergehen?
\end{teile}
\end{aufgabe}

\begin{aufgabe}{Ich kann zwei Sparangebote mit linearem und exponentiellem Wachstum vergleichen.}
Konto A: 1000\,€, jedes Jahr kommen 60\,€ dazu. \quad Konto B: 1000\,€, jedes Jahr kommen 5\,\% Zinsen dazu, die mitverzinst werden.
\begin{teile}
\teil Berechne beide Kontostände nach 1, 2 und 3 Jahren.
\teil Rechne weiter: Nach welchem Jahr hat Konto B zum ersten Mal mehr Geld als Konto A?
\teil Welches Konto ist besser, wenn du 5 Jahre sparst, welches bei 10 Jahren? Begründe mit den Kontoständen.
\end{teile}
\end{aufgabe}
